\subsubsection{The inertial quadratic model}
\label{sec:momentum-filtering}

The implicit solve \cref{eq:abstract-momentum-step} has a quadratic
optimization interpretation.

\begin{proposition}
\label{prop:inertial-model}
The increment \(\Delta\theta^{n+1}\) solving
\cref{eq:abstract-momentum-step} is the unique minimizer of
\begin{equation}
  \mathcal Q(s)
  =\frac{q^n}{Q^n}(g^n)^{\mathsf T}s
  +\frac12 s^{\mathsf T}B_\alpha^ns
  +\frac{1}{2\eta}\norm{s-\beta M^{n+1}p^n}_{(M^{n+1})^{-1}}^2 .
  \label{eq:inertial-quadratic-model}
\end{equation}
Appending the momentum after the solve gives instead
\begin{equation}
  \Delta_{\mathrm{app}}^{n+1}
  =-\eta\frac{q^n}{Q^n}
    \bigl((M^{n+1})^{-1}+\eta B_\alpha^n\bigr)^{-1}g^n
   +\beta M^{n+1}p^n,
  \label{eq:frozen-appended-comparator}
\end{equation}
and the two increments satisfy
\begin{align}
  \Delta_{\mathrm{app}}^{n+1}-\Delta\theta^{n+1}
  &=\beta\eta\bigl((M^{n+1})^{-1}+\eta B_\alpha^n\bigr)^{-1}
    B_\alpha^nM^{n+1}p^n,
  \label{eq:appended-model-step-defect}\\
  \mathcal Q(\Delta_{\mathrm{app}}^{n+1})-\mathcal Q(\Delta\theta^{n+1})
  &=\frac{1}{2\eta}
  \norm{\Delta_{\mathrm{app}}^{n+1}-\Delta\theta^{n+1}}
       _{(M^{n+1})^{-1}+\eta B_\alpha^n}^2\geq0.
  \label{eq:appended-model-excess}
\end{align}
\end{proposition}

\begin{proof}
The Hessian of \(\mathcal Q\) is
\(B_\alpha^n+\eta^{-1}(M^{n+1})^{-1}\succ0\), and its optimality equation is
\cref{eq:abstract-momentum-step}. Subtracting the two increments and using
\(\bigl((M^{n+1})^{-1}+\eta B_\alpha^n\bigr)M^{n+1}
 =I_d+\eta B_\alpha^nM^{n+1}\) gives
\cref{eq:appended-model-step-defect}. Expanding \(\mathcal Q\) about its
minimizer gives \cref{eq:appended-model-excess}.
\end{proof}

\begin{remark}
\label{rem:inertial-gauss-newton}
For a fixed \(\alpha\), let \(\Psi_\alpha(\theta)\) stack
\(\sqrt{2(1-\alpha)}\,Q(\theta)\) on \(\sqrt{\alpha}\,\psi(\theta)\), so that
\(\tfrac12\norm{\Psi_\alpha}^2=F+C\), its Gauss--Newton matrix is
\(B_\alpha\), and its gradient is \(g\). At exact
tracking \(q^n=Q^n\), the shifted model \(\mathcal Q(s)+F+C\) is a
Gauss--Newton model for the residual map \(\Psi_\alpha\), together with a
proximal term centered at \(\beta M^{n+1}p^n\) and measured in
\(\norm{\cdot}_{(M^{n+1})^{-1}}\); otherwise the linear term is scaled by
\(q^n/Q^n\). Levenberg--Marquardt methods for neural networks also combine curvature with
adaptive momentum \cite{pooladzandi2022improving}, without the scalar
coupling that produces the energy law below.
\end{remark}

To see how the solve acts on memory, let \(b_j\geq0\) be the eigenvalues of
\((M^{n+1})^{1/2}B_\alpha^n(M^{n+1})^{1/2}\) and let \(u_j\) be the
coordinates of \((M^{n+1})^{1/2}p^n\) in a corresponding orthonormal
eigenbasis. In those coordinates the solve multiplies the \(j\)th component
of the stored momentum by \(\beta/(1+\eta b_j)\), and
\begin{equation}
  \mathcal Q(\Delta_{\mathrm{app}}^{n+1})-\mathcal Q(\Delta\theta^{n+1})
  =\frac{\beta^2\eta}{2}
  \sum_j\frac{b_j^2}{1+\eta b_j}u_j^2 .
  \label{eq:spectral-memory-coefficient}
\end{equation}

\begin{remark}
\label{rem:memory-damping}
Directions with larger \(b_j\) damp the stored momentum more strongly, and
where \(b_j=0\) the coefficient remains \(\beta\). The excess in
\cref{eq:spectral-memory-coefficient} is measured in the quadratic model, not
in the modified energy of \cref{thm:monotone-mobility-dissipation}.
\end{remark}
